\section{Splitting along the crossing locus}\label{geometry}
Let $k$ be any field of characteristic $p>0$. A \emph{splinter} is a ring over which every module-finite extension splits.
\subsection{Trace ideals and coefficient extraction}
Trace ideals are uniformly compatible with Cartier operators \cite{DT2023,PS2023}. The following formula works over an arbitrary coefficient field.
\begin{lemma}\label{cartier-trace}
Let $q=p^e$.  Every additive map $C:R\to R$ satisfying
$C(a^qb)=aC(b)$ preserves $J(S/R)$.
\end{lemma}
\begin{proof}
If $c=\phi(1)$, the map
\[
 S\longrightarrow R,\qquad s\longmapsto C\bigl(\phi(s^q)\bigr)
\]
is $R$-linear and takes $1$ to $C(c)$.  Additivity uses commutativity of
$S$, and scalar compatibility follows from
$\phi((as)^q)=a^q\phi(s^q)$.
\end{proof}

For $P=k[[X_1,\ldots,X_d]]$, choose a multi-index $\gamma$ with
$0\le\gamma_i<q$.  Choose also a $k^q$-linear functional $k\to k^q$ and
write $\lambda$ for its composite with inverse Frobenius $k^q\to k$.  Then
\[
 L_\gamma\!\left(\sum_\nu b_\nu X^\nu\right)
   =\sum_{\delta\in\mathbf N^d}
       \lambda(b_{q\delta+\gamma})X^\delta
\]
is a Cartier operator because coefficient convolution gives
$L_\gamma(A^qB)=A L_\gamma(B)$.
For $R=P/(f)$ it induces the operator
\[
 \overline H\longmapsto
 \overline{L_\gamma(f^{q-1}H)}.
\]
The map is well-defined because replacing $H$ by $H+fA$ changes the
output by $fL_\gamma(A)$.  This is the coefficient calculation underlying
Fedder's criterion \cite{Fed}.  If the coefficient of $X^\gamma$ in the
input is a nonzero element $b\in k$, extend $b$ to a $k^q$-basis of $k$
and choose the functional to send $b$ to $1$.  The extracted series then
has constant term $1$.  No assumption on $[k:k^p]$ is involved.

\begin{lemma}\label{conductor-lower}
Let $R$ be reduced and $R\subseteq B\subseteq Q(R)$ a module-finite birational extension. If $B$ is a splinter, then for every module-finite extension $R\subseteq S$ one has
\[
 (R:B)=J(B/R)\subseteq J(S/R).
\]
A subring which is a module direct summand of a reduced splinter is itself a splinter.
\end{lemma}
\begin{proof}
A functional on $B$ becomes scalar multiplication on $Q(R)$, so its values at $1$ form $(R:B)$. In $S\otimes_R Q(R)$, take the subring generated by $B$ and the image of $S$. It is module-finite over $B$ and hence has a $B$-linear retraction to $B$. Restrict to $S$ and multiply by any element of $(R:B)$ to obtain the desired functional $S\to R$.

For the second assertion let $B\subseteq H$ split as $B$-modules with $H$ a reduced splinter. For a module-finite extension $B\subseteq T$, the map $\Spec T\to\Spec B$ is surjective by lying over. It stays surjective after base change to $H$, so $H\to H\otimes_B T$ has nil kernel and is injective. Compose $T\to H\otimes_B T$ with the retractions to $H$ and to $B$.
\end{proof}

\subsection{Coordinate crossings and suspensions}
\begin{theorem}\label{cross-split}
Let $s\ge2$ and let
\[
 R=k[[x_1,\ldots,x_s,z_1,\ldots,z_n]]/(x_ix_j\mid i<j),
 \qquad P_0=(x_1,\ldots,x_s)R .
\]
Every evaluation ideal on $R$ is either $R$ or $P_0$, and the normalization realizes $P_0$. Splitting is detected at $P_0$, and $R$ has $\K$.
If $0\ne g\in(z_1,\ldots,z_n)^2$, the hypersurface
$k[[x,y,z_1,\ldots,z_n]]/(xy-g)$ is a domain and a splinter.
\end{theorem}
\begin{proof}[Proof for coordinate crossings]
Write $C=k[[z_1,\ldots,z_n]]$.
Represent $R$ as the tuples in $\prod_i C[[x_i]]$ having a common value
at $x_i=0$.  If $a\in J(S/R)\setminus P_0$, this common value has a
nonzero coefficient at some $z^\alpha$.  Choose $q>\max_j\alpha_j$
(any $q>1$ when $n=0$).  On each branch extract the residue exponents
$(0,\alpha)$ in $(x_i,z)$, using the same $\lambda$ on the coefficients.
The value at $x_i=0$ of each output depends only on the common input in
$C$.  The branch operators therefore preserve $R$ and define a Cartier
operator carrying $a$ to a unit.  Lemma~\ref{cartier-trace} implies
\[
 J(S/R)\ne R\quad\Longrightarrow\quad J(S/R)\subseteq P_0.
\]
The normalization is $B=\prod_i C[[x_i]]$. Indeed $B$ is a product of normal domains, is module-finite over $R$ by adjoining its coordinate idempotents, and agrees with $R$ after inverting $(x_i)_i$. Its conductor is $P_0$, since multiplying a tuple by each coordinate idempotent shows that its common value must vanish for the product to stay in $R$. Each factor is a regular splinter \cite{Andre}, so $B$ is a splinter by applying retractions factor by factor. Lemma~\ref{conductor-lower} yields $P_0\subseteq J(S/R)$ and $J(B/R)=P_0$. This proves the dichotomy.
Since evaluation ideals commute with localization, the extension splits
over $R$ if and only if it splits over $R_{P_0}$.

It remains to apply the curve theorem when $\pd_R S<\infty$.
Write $K=\operatorname{Frac}(C)$.  For $j\ge1$, the quotient $R/P_0^j$
is finite free over $C$, with basis $1$ and the pure powers $x_i^t$ for
$1\le t<j$.  Tensoring with $K$ gives the local Artinian truncations of
the union of axes over $K$.  Taking inverse limits gives
\[
 \widehat{R_{P_0}}\simeq\prod_K K[[x_i]].
\]
Localization and completion preserve the module-finite extension and its finite
projective dimension.  Theorem~\ref{branch-split} gives a retraction over this completed
curve.  Faithfully flat descent gives one over $R_{P_0}$, and the
localization criterion gives one over $R$.
\end{proof}

The strong $F$-regularity of the suspension $xy-g(z)$ is familiar in the
usual geometric setting \cite[p.~5]{HP2017}.  We retain the direct argument below to
cover arbitrary coefficient fields and all finite targets allowed in the
definition of $\K$.

\begin{proof}[Proof for suspensions]
Let $f=xy-g(z)$ with $0\ne g\in(z)^2$.  Replacing $xy$ by $g$ preserves
or increases total degree and gives a convergent normal form
\[
 c=\sum_{a\ge0}x^a A_a(z)+\sum_{b\ge1}y^b B_b(z).
\]
For uniqueness, substitute $y=g/x$ in $k((x))[[z]]$.  At each fixed
$z$-degree only finitely many negative powers of $x$ contribute, because
$g$ has positive order.  The coefficient of $x^{-b}$ is $g^bB_b$, while
that of $x^a$ is $A_a$.  Since $g\ne0$, the normal form is unique and the
quotient embeds in the domain $k((x))[[z]]$.

Let $c\ne0$ and suppose $A_a\ne0$.  Choose $\beta$ such that
$[z^\beta]g^aA_a\ne0$, and take $q$ larger than $a$ and all $\beta_i$.
The coefficient of $x^{q-1}y^{q-1-a}z^\beta$ in $f^{q-1}c$ is
\[
 (-1)^a\binom{q-1}{a}[z^\beta]g^aA_a
   =[z^\beta]g^aA_a\ne0.
\]
Indeed, the difference of the $x$- and $y$-exponents selects
$x^aA_a$, and
$(1+T)^{q-1}=\sum_{i=0}^{q-1}(-T)^i$ in characteristic $p$.
If only the second sum is nonzero, interchange $x$ and $y$.
Coefficient extraction therefore sends every nonzero element of $R$
to a unit by a suitable Cartier operator.

For any module-finite extension over a domain, $J(S/R)$ contains a nonzero
element: choose a linear functional on $S\otimes_R\operatorname{Frac}R$
taking $1$ to $1$ and clear its denominators on a finite generating set
of $S$.  Cartier compatibility now gives $J(S/R)=R$.
\end{proof}
\begin{corollary}\label{quadratic-rank}
Let $k$ have odd characteristic and write $X=(X_1,\ldots,X_d)$. If $f\in(X)^2$ has quadratic part of rank at least two, the ring $k[[X]]/(f)$ has $\K$.
It is a splinter if the quadratic rank is at least three, or if the residual series after splitting off a hyperbolic plane over a finite separable coefficient extension is nonzero.
\end{corollary}
\begin{proof}
In odd characteristic, diagonalize the quadratic form.  After a separable
extension of degree at most two, a nondegenerate binary summand becomes
hyperbolic.  Write the quadratic part as $uv+Q(z)$.
At each degree $m\ge3$, write the terms involving $u$ or $v$ as
$uA+vB$, with $A,B$ homogeneous of degree $m-1$.  The coordinate changes
$v\mapsto v-A$ and $u\mapsto u-B$ remove these terms without changing
the lower-degree terms.  Their successive compositions converge in the
maximal-ideal topology to a formal automorphism.  The equation becomes
$uv-g(z)$, with quadratic part of $g$ equal to $-Q$.
The crossing case applies if $g=0$ and the suspension case applies if $g\ne0$.
Quadratic rank at least three forces $Q\ne0$.  Finally, the finite
extension of coefficient fields commutes with power series and is
faithfully flat, so splitting descends by flat base change for the evaluation ideal.
\end{proof}

\subsection{Pinch points and the conductor}
\begin{theorem}\label{pinch}
Let $C=k[[z_1,\ldots,z_s,w_1,\ldots,w_n]]$, $s\ge1$, and $u\in C^\times$. Put $g=u z_1\cdots z_s$ and $R=C[[x,y]]/(x^2-y^2g)$.
If $p$ is odd, every evaluation ideal on $R$ is either $R$ or $(x,y)$. Its normalization realizes $(x,y)$, splitting is detected at $(x,y)$, and $R$ has $\K$. In characteristic two each of these hypersurfaces admits a nonsplit module-finite extension of generic rank four and projective dimension one.
\end{theorem}
\begin{proof}
First work in odd characteristic, and put
$C=k[[z_1,\ldots,z_n]]$.  More generally, let $0\ne g\in C$ satisfy
\begin{equation}\label{branchcond}
 \text{for each }0\ne c\in C\text{ some }q=p^e\text{ satisfies }
 c g^{(q-1)/2}\notin(z_1^q,\ldots,z_n^q).
\end{equation}
Write $R=C[[x,y]]/(x^2-y^2g)$ and $P_0=(x,y)R$.
Division by the monic equation in $x$ represents any element outside
$P_0$ as
$H=c+yA(y,z)+xB(y,z)$ with $0\ne c\in C$.
Choose $q$ as in \eqref{branchcond}, set $h=(q-1)/2$, and choose a
nonzero coefficient of $cg^h$ at $z^\beta$ with all $\beta_i<q$.
In $(x^2-y^2g)^{q-1}H$, the coefficient at
$x^{q-1}y^{q-1}z^\beta$ is precisely this coefficient.
The $xB$ term has odd $x$-degree, the $yA$ term overshoots the required
$y$-degree once the $x$-degree is fixed, and the contributing binomial
coefficient is
$(-1)^h\binom{q-1}{h}=1$.
Extraction sends $H$ to a unit.  Hence every proper $J(S/R)$ is
contained in $P_0$, and splitting is detected at $P_0$.

With $K=\operatorname{Frac}C$, truncation in $x,y$ gives
\[
 \widehat{R_{P_0}}\simeq K[[x,y]]/(x^2-gy^2).
\]
After the separable extension adjoining $\sqrt g$, if needed, the
coordinates $x+\sqrt g\,y$ and $x-\sqrt g\,y$ identify this ring with
a node.  Theorem~\ref{branch-split} therefore applies to every module-finite extension of finite
projective dimension.  Descent and the preceding detection criterion
give $\K$ for $R$.

For $g=u z_1\cdots z_s$ with $u\in C^\times$, choose a nonzero term
of $c$ of least total degree.  Its coefficient in $cu^h$ is multiplied
by the nonzero constant term of $u^h$, since $c$ has no terms of smaller
degree.  Multiplication by $(z_1\cdots z_s)^h$ translates its first
$s$ exponents by $h$.  All translated exponents are less than $q$
for sufficiently large $q=p^e$.  This proves \eqref{branchcond}, including the case of additional smooth variables.

For the ideal calculation, the coordinate change $z_1\mapsto u z_1$ absorbs the unit. Set $A=C[[y]]$ and $B=A[t]/(t^2-z_1\cdots z_s)$. The substitutions $z_i=v_i^2$ and $t=v_1\cdots v_s$ identify $B$ with the subring of
$H=k[[v_1,\ldots,v_s,y,w_1,\ldots,w_n]]$ consisting of series whose $v_i$-exponents are either all even or all odd. The two parity classes give the summands in $B=A\oplus At$. Projection onto these monomials is $B$-linear, since their parity vectors form a subgroup of $(\mathbf Z/2)^s$. Hence $B$ is a splinter by Lemma~\ref{conductor-lower}. For odd $p$, $B$ is the fixed ring of the sign changes $v_i\mapsto\epsilon_i v_i$ with $\prod_i\epsilon_i=1$. Indeed these fix precisely the two stated parity classes. Thus $B$ is normal, since $H$ is normal and an element of $\operatorname{Frac}B$ is fixed by every sign change.

The substitution $x=yt$ identifies $R=A\oplus yAt$ inside $B$. It is birational since $t=x/y$, and its conductor is $yB=(x,y)R$. To see this, $a+ybt\in R$, $a,b\in A$, belongs to the conductor exactly when its product with $t$ lies in $R$, or $a\in yA$. Thus $B$ is the normalization for odd $p$. Lemma~\ref{conductor-lower} supplies the lower bound $(x,y)\subseteq J(S/R)$, and the Cartier calculation gives the upper bound for every proper evaluation ideal. In particular $J(B/R)=(x,y)$, so $\pd_R B=\infty$ since $R$ has $\K$.

In characteristic two the same parity description and $R=A\oplus yAt$ give $t=x/y\notin R$ and $t^2\in R$. Apply Theorem~\ref{realize}. For the nonreduced ring $k[[x,y]]/(y^2)$, apply
Lemma~\ref{model} directly with $a=x$ and $\eta=y\notin xR$.
The resulting rank-two algebra has a minimal differential of size
$5\times3$, and the relation $x^6(y+xu)=0$ rules out a retraction.
\end{proof}
